\documentclass[11pt]{article}
\usepackage[margin=1in]{geometry}
\usepackage{amsmath,amssymb,amsthm,mathtools}
\usepackage{microtype}
\usepackage[hidelinks]{hyperref}
\newtheorem{theorem}{Theorem}
\newtheorem{proposition}{Proposition}
\newtheorem{remark}{Remark}
\newcommand{\R}{\mathbb R}
\newcommand{\D}{\mathbb D}
\newcommand{\Fix}{\operatorname{Fix}}
\newcommand{\distH}{d_{\mathrm H}}
\newcommand{\cH}{\mathcal H}
\newcommand{\cR}{\mathbf R}
\newcommand{\cT}{\mathbf T}
\title{Connected continuation of the smooth gCLM profiles\\
and a locally unique analytic branch at the CLM anchor}
\author{Discovery note for arXiv:2305.05895}
\date{17 August 2026}

\begin{document}
\maketitle

\begin{abstract}
The source paper constructs, for every $a\leq1$, a fixed point of a
continuous map $\cR_a$ on one compact convex set $\D$, and conjectures
uniqueness, pointwise monotonicity in $a$, and a continuous choice of fixed
points.  We prove two partial results.  First, on every compact parameter
interval the fixed-point graph has a connected component projecting onto the
whole interval; on $[0,1]$ it joins the unique CLM and De Gregorio anchors,
and every selection is continuous at both anchors.  Second, the exact CLM
spectrum recently computed in arXiv:2607.19762 makes the scale-fixed profile
equation analytically invertible, giving a locally unique real-analytic
branch of strongly regular normalized profiles near $a=0$.  We also explain
why neither conclusion yet gives the source's global monotone graph.
\end{abstract}

\section{Source conjecture and result}

Let $\rho(x)=(1+|x|)^{-1/2}$, let $V$ be the source paper's Banach space with
norm $\|f\|_\rho=\|\rho f\|_\infty$, and let $\D\subset V$ be its compact
convex invariant set.  The map $\cR_a:\D\to\D$ is defined in Section 3 of
the source for every $a\leq1$.  Put
\[
 \mathcal F_a=\Fix(\cR_a),\qquad
 \mathcal G_I=\{(a,f)\in I\times\D:f=\cR_a(f)\}
\]
for a compact interval $I\subset(-\infty,1]$.
Conjecture 3.12 asks for pointwise parameter monotonicity of all these fixed
points (hence uniqueness of every $\mathcal F_a$) and for a continuous
selection $a\mapsto f_a$.

\begin{theorem}[Unconditional continuation and anchor localization]
\label{thm:continuation}
For every compact interval $I=[A,B]\subset(-\infty,1]$:
\begin{enumerate}
\item $\mathcal G_I$ is nonempty and compact, and it has a compact connected
component $C_I$ whose projection onto the parameter coordinate is all of
$I$.
\item On $I=[0,1]$, such a component meets $(0,f_0)$ and $(1,f_1)$, where
$f_0$ and $f_1$ are the unique fixed points at the two endpoints.
\item In the Hausdorff metric of $V$,
\[
 \distH(\mathcal F_a,\{f_0\})\longrightarrow0\quad(a\to0),
 \qquad
 \distH(\mathcal F_a,\{f_1\})\longrightarrow0\quad(a\to1^-).
\]
Consequently every arbitrary choice $f_a\in\mathcal F_a$ is continuous at
$a=0$ and at $a=1$.
\end{enumerate}
Moreover the conjectured order is exactly true for the two explicit anchors
$a=0$ and $a=\tfrac12$:
\[
 f_0(x)=\frac1{1+x^2}\ \geq\
 f_{1/2}(x)=\frac4{(2+x^2)^2}\qquad(x\in\R).
\]
\end{theorem}

The second result concerns the equivalent profile equation in the standard
normalization with time exponent $-1$.  Let $X=H^2_{\rm odd}(\R)$ with the
equivalent norm $\|\phi\|_X^2=\|\phi\|_2^2+\|\phi''\|_2^2$, and let
\[
 E=\{\Omega\in X:y\Omega'\in X\}
\]
with its graph norm.  For odd $\Omega$, define
$U_\Omega'=\cH\Omega$ and $U_\Omega(0)=0$, where $\cH$ is the Hilbert
transform, and set
\[
 \Phi(a,\Omega,c)=
 \Omega+(cy+aU_\Omega)\Omega'-\Omega\cH\Omega.
\]
The exact CLM point is
\[
 (a,\Omega,c)=(0,\Omega_0,1),\qquad
 \Omega_0(y)=-\frac{y}{y^2+\tfrac14}.
\]

\begin{theorem}[Locally unique analytic CLM branch]
\label{thm:ift}
There are $\varepsilon>0$, a graph-norm neighborhood $\mathcal U$ of
$\Omega_0$ in $E$, and a neighborhood $J$ of $1$ such that, for every
$|a|<\varepsilon$, the system
\[
 \Phi(a,\Omega,c)=0,
 \qquad \Omega'(0)=\Omega_0'(0)=-4
\]
has exactly one solution $(\Omega_a,c_a)\in\mathcal U\times J$.  The map
$a\mapsto(\Omega_a,c_a)$ is real analytic.
\end{theorem}

This is local uniqueness in the strong profile graph norm.  It does not
assert that every Schauder fixed point in $\D$ lies in this neighborhood.

\section{Idea of the proof}

The source already supplies more than separate Schauder fixed points.  Its
compactness lemma makes $\D$ a compact convex set, and the proof of its
uniform tail lemma records joint continuity of $(a,f)\mapsto\cR_a(f)$,
including the removable formula at $a=0$.  Browder's parametric strengthening
of Schauder then forces one connected component of fixed points to cross any
compact parameter interval.  Uniqueness at an endpoint upgrades compactness
to convergence of the entire fixed-point fiber there.

For the strong local statement, the profile residual is analytic on a graph
domain.  The exact spectral theorem at $a=0$ says that the only zero mode of
its linearization is spatial scaling and that this mode is semisimple.  The
unknown speed $c$ supplies that missing spectral direction, while the trace
$\Omega'(0)=-4$ fixes the scale.  The augmented derivative is therefore an
isomorphism, so the analytic implicit-function theorem applies.

\section{Proof of Theorem \ref{thm:continuation}}

The source proves continuity of $f\mapsto\cR_a(f)$ on $\D$.  More
importantly, in the proof of its uniform-tail Lemma 4.7 it observes that if
$a_n\to a$ and $f_n\to f$ in $V$, then
\[
 \cR_{a_n}(f_n)\longrightarrow\cR_a(f)\quad\hbox{in }V.
\]
At $a=0$ this uses continuity of $(1-t)^{1/t}$ at zero.  The estimates in
that proof are uniform when $a$ stays in a compact subinterval of
$(-\infty,1]$.  Hence $(a,f)\mapsto\cR_a(f)$ is jointly continuous on
$I\times\D$.

The parameter space $I$ is compact and connected and $\D$ is a nonempty
compact convex subset of the locally convex space $V$.  Browder's parametric
fixed-point theorem therefore says that $\mathcal G_I$ has a connected
component whose projection is $I$.  Joint continuity also makes
$\mathcal G_I$ closed in the compact set $I\times\D$, hence compact.  This
proves the first assertion.  For $I=[0,1]$, the source proves that the endpoint
fibers are the singletons $\{f_0\}$ and $\{f_1\}$, so a component projecting
onto the whole interval must meet both indicated points.

For anchor localization, suppose for example that the first Hausdorff limit
fails.  There are $a_n\to0$ and $f_n\in\mathcal F_{a_n}$ with
$\|f_n-f_0\|_\rho\geq\delta>0$.  Compactness of $\D$ gives a subsequence
$f_n\to f_*$.  Joint continuity and $f_n=\cR_{a_n}(f_n)$ imply
$f_*=\cR_0(f_*)$, so endpoint uniqueness gives $f_*=f_0$, a contradiction.
The proof at $a=1$ is identical.

Finally, with $t=x^2$,
\[
 \frac1{1+x^2}-\frac4{(2+x^2)^2}
 =\frac{(t+2)^2-4(t+1)}{(1+t)(2+t)^2}
 =\frac{x^4}{(1+x^2)(2+x^2)^2}\geq0.
\]

\section{Proof of Theorem \ref{thm:ift}}

First we verify the nonlinear mapping.  For $\Omega\in X$, put
\[
 Q_\Omega(y)=\frac{U_\Omega(y)}y
 =\int_0^1(\cH\Omega)(ty)\,dt,
\]
with the evident value at $y=0$.  For $k=0,1,2$, dilation and Minkowski give
\[
 \|\partial_y^k Q_\Omega\|_2
 \leq \|\partial_y^k\cH\Omega\|_2
       \int_0^1t^{k-1/2}\,dt.
\]
Thus $\Omega\mapsto Q_\Omega$ is bounded on $H^2$.  Since $H^2(\R)$ is a
Banach algebra and $U_\Omega\Omega'=Q_\Omega(y\Omega')$, the map
\[
 \Phi:\R\times E\times\R\longrightarrow X
\]
is real analytic.

Let $L_0$ be the closed CLM linearization on $X$ with domain $E$.  Direct
differentiation gives
\[
 D_\Omega\Phi(0,\Omega_0,1)=-L_0,
 \qquad \partial_c\Phi(0,\Omega_0,1)=s:=y\Omega_0'.
\]
The exact spectral results of arXiv:2607.19762 show that zero is an isolated
semisimple eigenvalue of $L_0$, its eigenspace is one-dimensional and spanned
by $s$, and the restriction of $L_0$ to the complementary Riesz subspace is
boundedly invertible (from $X$ into the graph domain).  Let $P$ be the Riesz
projection at zero and write
\[
 X=\operatorname{span}\{s\}\oplus X_1,
 \qquad E=\operatorname{span}\{s\}\oplus E_1.
\]

The derivative of the augmented map
$(\Phi,\Omega'(0)+4)$ is
\[
 \mathcal A(\phi,d)=(-L_0\phi+ds,\phi'(0)).
\]
It is an isomorphism.  Indeed, given $(g,r)\in X\times\R$, decompose
$g=\beta s+g_1$ with $g_1\in X_1$, take
\[
 d=\beta,
 \qquad \phi_1=-\bigl(L_0|_{E_1}\bigr)^{-1}g_1,
 \qquad \phi=\alpha s+\phi_1,
\]
and choose $\alpha$ so that $\phi'(0)=r$.  This is possible and unique since
\[
 s'(0)=\Omega_0'(0)=-4\neq0.
\]
All operations are bounded; uniqueness follows from the same decomposition.
The analytic implicit-function theorem now proves the claim.

\section{Why this is not yet the full conjecture}

The kernel of the source's linear transform is
\[
 K(x,y)=\frac yx\log\left|\frac{x+y}{x-y}\right|-2.
\]
For $t=y/x$, it tends to $-2$ as $t\to0$ but equals
$2/(3t^2)+O(t^{-4})>0$ as $t\to\infty$.  Thus $\cT$ is not positivity
preserving, blocking a direct order iteration for the pointwise conjecture.
The exact linearized inverse at $a=0$ is strongly nonnormal and has no known
positivity property either.

Nor does a connected fixed-point continuum automatically contain a
continuous section.  On $[0,1]\times[0,1]$, set
\[
 g(t)=\frac{1+\sin(1/t)}2\quad(t>0),\qquad
 F(t,x)=(1-t)x+tg(t),\qquad F(0,x)=x.
\]
This is jointly continuous.  Its fixed set is
\[
 \{(t,g(t)):t>0\}\cup(\{0\}\times[0,1]),
\]
the topologist's sine curve: it is connected and projects onto the whole
parameter interval, but it has no continuous section.  A proof of the full
source conjecture therefore needs additional analytic sign or strong
compactness structure.

\section*{References}

\begin{enumerate}
\item D. Huang, X. Qin, X. Wang, and D. Wei,
\emph{Self-similar finite-time blowups with smooth profiles of the generalized
Constantin--Lax--Majda model}, Arch. Ration. Mech. Anal. 248 (2024), article
22; arXiv:2305.05895.  Conjecture 3.12 and Sections 3--5.
\item E. Solan and O. N. Solan, \emph{Browder's Theorem: from
One-Dimensional Parameter Space to General Parameter Space},
arXiv:2210.16369.  Parametric fixed-point theorem for compact connected
parameter spaces and compact convex subsets of locally convex spaces.
\item J. Xu, \emph{The spectral picture of self-similar collapse in the
Constantin--Lax--Majda equation}, arXiv:2607.19762 (2026).  Exact
origin-$H^2$ spectrum, simple symmetry modes, and resolvent at $a=0$.
\end{enumerate}

\end{document}
